\documentclass[10pt,a4paper]{amsart}
\usepackage[margin=19mm]{geometry}
\usepackage{amsmath,amssymb,mathtools}
\usepackage[scr=rsfs,cal=euler]{mathalfa}
\usepackage{xfrac}
\usepackage[unicode,hidelinks,bookmarks=false]{hyperref}
\newcommand{\ie}{i.e.\ }
\newcommand{\cf}{cf.\ }
\newcommand{\resp}{respectively\ }
\newcommand{\Subsection}{\subsection}
\providecommand{\nobreakdash}{\nobreak\hbox{-}}
\setlength{\emergencystretch}{2em}
\multlinegap0pt
\usepackage{amssymb}
\usepackage[shortlabels]{enumitem}
\usepackage{mathtools}
\usepackage{tikz}
\usepackage{booktabs}
\usepackage{siunitx}
\newtheorem{theorem}{Theorem}
\newtheorem{proposition}{Proposition}
\usepackage{cleveref}
\crefname{theorem}{Theorem}{Theorems}
\crefname{proposition}{Proposition}{Propositions}
\newcommand{\E}{\mathrm{E}}
\newcommand{\N}{\mathbb{N}}
\newcommand{\Eq}[1]{\mathrm{E}#1}
\newcommand{\formaleq}{\simeq}
\newcommand{\nmodels}{\not\models}
\newcommand{\op}{\diamond}
\newcommand{\w}{\mathrm{w}}
\newcommand{\x}{\mathrm{x}}
\newcommand{\y}{\mathrm{y}}
\newcommand{\z}{\mathrm{z}}
\title{The law E73 does not imply cube associativity E4380}
\author[ETP contributors]{Matthew Bolan, Joachim Breitner, Jose Brox, Nicholas Carlini, Mario Carneiro, Floris van Doorn, Martin Dvorak, Andr\'es Goens, Aaron Hill, Harald Husum, Hern\'an Ibarra Mejia, Zoltan A. Kocsis, Bruno Le Floch, Amir Livne Bar-on, Lorenzo Luccioli, Douglas McNeil, Alex Meiburg, Pietro Monticone, Pace P. Nielsen, Emmanuel Osalotioman Osazuwa, Giovanni Paolini, Marco Petracci, Bernhard Reinke, David Renshaw, Marcus Rossel, Cody Roux, J\'er\'emy Scanvic, Shreyas Srinivas, Anand Rao Tadipatri, Terence Tao, Vlad Tsyrklevich, Fernando Vaquerizo-Villar, Daniel Weber, Fan Zheng}
\date{Source: arXiv:2512.07087v2 (CC BY 4.0)}
\newcommand{\modelsfin}{\models_{\mathrm{fin}}}
\begin{document}
\maketitle

\noindent\emph{Source and licence.} The Equational Theories Project: Advancing Collaborative Mathematical Research at Scale. Version arXiv:2512.07087v2 (CC BY 4.0), https://arxiv.org/abs/2512.07087v2. The exact-version arXiv abstract page was retrieved and its CC BY 4.0 article licence was independently checked on 2026-10-10. Licence and source notices: \texttt{licenses/arxiv-2512.07087-CC-BY-4.0.txt}.\par\smallskip

\noindent\textit{Editorial scope.} The counted unit is the original Proposition 73-4380 (source0.tex lines 784--785), with its complete author proof at lines 787--811. The statement, proof and retained source passages are copied line for line without mathematical rewriting. Uncounted source-local context is retained at 749--772, 1778, 1817, 819--819. The source article remains the work cited above; the title of this delivery describes the selected result. The unchanged native source is bundled as \texttt{sources/190548/source0.tex}. The complete abstract greedy-algorithm proof is included. Source line 819 describes these as human-written greedy arguments verified with ATPs and Lean; this statement concerns this argument and is not a claim that the entire collaborative paper had no AI assistance. The Skolem reference is transcribed from reference 63 of the same fixed-version arXiv HTML. The uncounted all-equations table of r002 is replaced by the two verbatim source rows defining E73 and E4380, typeset in a small editorial align environment. All other table laws are irrelevant to this counted proof. Running headers use a short contributor name while the title page retains all source authors.
We have found \emph{greedy extension methods}, or \emph{greedy methods} for short, to be a powerful way to refute implications, especially when the carrier $M$ is allowed to be infinite.  Such constructions have a long history in model theory, with possibly the earliest\footnote{We thank Stanley Burris for this reference.} such construction due to Skolem \cite{skolem}. A basic implementation of this method is as follows.  To build a magma operation $\op \colon M \times M \to M$ that satisfies one law $\E$ but not another ${\E'}$, one can first consider \emph{partial magma operations} $\op \colon \Omega \to M$, defined on some subset $\Omega$ of $M \times M$. Thus $x \op y$ is defined if and only if $(x,y) \in \Omega$. A magma operation is then simply a partial operation which is \emph{total} in the sense that $\Omega = M \times M$.  We say that a partial magma operation is \emph{finitely supported} if $\Omega$ is finite.

In the language of first-order logic (in which functions and relations must be total), it is convenient to view a partial magma operation as a ternary relation $R(x,y,z)$ on $M$ with the axiom that $R(x,y,z) \wedge R(x,y,z') \implies z=z'$ for all $x,y,z \in M$.  The support $\Omega$ is then the set of $(x,y)$ for which $R(x,y,z)$ holds for some (necessarily unique) $z$, which one can then take to be the definition of $z = x \op y$.

We say that one partial operation $\op' \colon \Omega' \to M$ \emph{extends} another $\op \colon \Omega \to M$ if $\Omega'$ contains $\Omega$, and $x \op y = x \op' y$ whenever $x \op y$ (and hence $x \op' y$) are defined. Given a sequence $\op_n \colon \Omega_n \to M$ of partial operations, each of which is an extension of the previous, we can define the \emph{direct limit} $\op_\infty \colon \bigcup_n \Omega_n \to M$ to be the partial operation defined by $x \op_\infty y \coloneqq x \op_n y$ whenever $(x,y) \in \Omega_n$.
\looseness=-1

Abstractly, the greedy algorithm strategy can now be described as follows.

\begin{theorem}[Abstract greedy algorithm]\label{greedy-abstract} Let $\E,{\E'}$ be equational laws, and let $\Gamma$ be a theory of first-order sentences regarding a partial magma operation $\op \colon \Omega \to M$ on a carrier $M$.  Assume the following axioms:
  \setlength{\parskip}{0pt}
\begin{itemize}
  \item[(i)] (Seed) There exists a finitely supported partial magma operation $\op_0 \colon \Omega_0 \to M$ satisfying $\Gamma$ that contradicts ${\E'}$, in the sense that there is some assignment of variables in ${\E'}$ in $M$ such that both sides of ${\E'}$ are defined using $\op_0$, but not equal to each other.
  \item[(ii)]  (Soundness)  If $\op_n \colon \Omega_n \to M$ is a sequence of partial magma operations satisfying $\Gamma$ with each $\op_{n+1}$ an extension of $\op_n$, and the direct limit $\op_\infty$ is total, then this limit satisfies $\E$.
  \item[(iii)] (Greedy extension)  If $\op \colon \Omega \to M$ is a finitely supported partial magma operation satisfying $\Gamma$, and $a,b \in M$, then there exists a finitely supported extension $\op' \colon \Omega' \to M'$ of $\op$ to a possibly larger carrier $M'$, and also satisfying $\Gamma$, such that $a \op' b$ is defined.
\end{itemize}
Then $\E \nmodels {\E'}$.
\end{theorem}

We remark that this greedy method seems to be inherently infinitary in nature, and does not seem well adapted to refute finite magma implications $\E \modelsfin {\E'}$.

\begin{proof}  We work on the countably infinite carrier $\N$.  By embedding the finitely supported operation $\op_0$ from axiom (i) into $\N$, we can assume without loss of generality that $\op_0$ has carrier $\N$.  By similar relabeling, we can assume in (iii) that $M' = M$ when $M=\N$, since any elements of $M' \setminus \N$ that
appear in $\Omega'$ can simply be reassigned to natural numbers that did not previously appear in $\Omega$.  We well-order the pairs in $\N \times \N$ by $(a_n,b_n)$ for $n=1,2,\dots$.  Iterating (iii) starting from $\op_0$, we can thus create a sequence of finitely supported magma operations $\op_0, \op_1, \dots$ on $\N$ satisfying $\Gamma$, with each $\op_{n+1}$ an extension of $\op_n$, and $a_n \op_n b_n$ defined for all $n \geq 1$.  Then the direct limit $\op_\infty$ of these operations is total, and does not satisfy ${\E'}$ thanks to axiom (i).  On the other hand, by axiom (ii) it satisfies $\E$, and the claim follows.
\end{proof}

\begin{align*}
    \x &\formaleq \y \op (\y \op (\x \op \y))  \label{eq73}\tag{E73} \\
        \x \op (\x \op \x) &\formaleq (\x \op \x) \op \x & \hbox{(Cube-associativity law)} \label{eq4380}\tag{E4380} \\
\end{align*}
We also found that human verification of the greedy extension property was a highly error-prone process, as the case analysis often included many delicate edge cases that were easy to overlook.  Both ATPs and the Lean formalization therefore played a crucial role in verifying the human-written greedy arguments, often revealing important gaps in those arguments that required either minor or major revisions to the rule set.

\begin{proposition}[$\Eq{73} \nmodels \Eq{4380}$]\label{73-4380} The law $\x \formaleq \y \op (\y \op (\x \op \y))$ \eqref{eq73} does not imply $\x \op (\x \op \x) \formaleq (\x \op \x) \op \x$ \eqref{eq4380}.
\end{proposition}

\begin{proof} To build a rule set $\Gamma$ that will imply $\Eq{73}$ when total, a natural first choice would be the single rule
\begin{itemize}
\item[1.] If $y \op (x \op y)$ is defined, then $y \op (y \op (x \op y))$ is defined and equal to $x$.
\end{itemize}
However, the greedy algorithm will fail just with this rule: if the partial operation has $x \op y$ and $z \op y$ both equal to some $w$ for some $x \neq z$, then any attempt to assign a value to $y \op (y \op w)$ will lead to a contradiction, as the above rule will force $y \op w$ to equal both $x$ and $z$.  Indeed, it is clear that $\Eq{73}$ forces all the right translation operators $R_y$ to be injective.  We therefore add this as an additional rule:
\begin{itemize}
\item[2.] If $x \op y$ and $z \op y$ are defined and equal, then $x=z$.
\end{itemize}
To avoid some unwanted edge cases, it is also convenient to impose the additional rule
\begin{itemize}
  \item[3.] If $x \op y$ is defined, it is not equal to $y$.
\end{itemize}
Unlike Rule 2, this rule is not forced by $\Eq{73}$, but can be enforced as part of the greedy construction.

The ruleset clearly satisfies the soundness axiom (ii) of \Cref{greedy-abstract}.  We now verify the greedy extension axiom (iii).  Let $\Omega,a,b$ be as in that axiom. We may assume that $a \op b$ is undefined, since we are done otherwise. We adjoin a new element $c$ to $M$ to create $M'$, and set $a \op' b = c$.  If we also have $b = d \op a$ for some $d$ (unique by Rule 2, and only possible for $a \neq b$ by Rule 3), set $a \op' c = d$ (this is necessary to retain Rule 1).  Of course, we also set $x \op' y = x \op y$ whenever $x \op y$ is already defined.

Since $c \not \in M$, it is clear that $\op'$ is a finitely supported partial magma operation on $M'$.  It is also clear that $\op'$ satisfies Rule 2 and Rule 3.   Now we case check Rule 1:
\begin{itemize}
\item Case 1: $x=c$ or $y=c$.  Not possible since no left multiplication with $c$ is defined.
\item Case 2: $x \op' y = c$.  Only possible when $x = a$, $y = b$, but then $y \op' (x \op' y)$ is undefined since $y = b \neq a$ if $d$ is defined.
\item Case 3. $y \op' (x \op' y) = c$.  Only possible when $y=a$ and $x=d$, and holds in this case.
\item Case 4: $x, y, x \op' y, y \op' (x \op' y) \neq c$: In this case $\op'=\op$ on all pairs, so the claim reduces to Rule 1 for $\op$, which holds by the induction hypothesis.
\end{itemize}
To conclude, we need to locate a seed $\op_0$ satisfying Rules 1,2,3 and containing a counterexample to $\Eq{4380}$.  One simple example is the carrier $\{0,1,2,3\}$ with $0 \op_0 0 = 1$, $0 \op_0 1 = 2$, $0 \op_0 2 = 0$, $1 \op_0 0 = 3$.
\end{proof}

\begin{thebibliography}{99}
\bibitem[63]{skolem} Th. Skolem. \emph{Logisch-kombinatorische Untersuchungen \"uber die Erf\"ullbarkeit oder Beweisbarkeit mathematischer S\"atze nebst einem Theoreme \"uber dichte Mengen}. Krist. Vid. Selsk. Skr. I, 1920, Nr. 4, 36 S. (1922)., 1922.
\end{thebibliography}
\end{document}
